%% kapa: exam=1 round=22 year=2011 period=2 form=A subject=appraisal_law subject_ko=부동산관계법규 seq=6 number=46
%% source: Q-Net 감정평가사 기출문제 (cst00309) artlSeq=5036790 file=extracted/1008/5036790/2011년 제22회 감정평가사 1차시험 기출문제/2교시.hwp
%% conversion: kordoc 4.24.51 -> tools/gen1.py (auto); review=proofread-gpt-astra; flags=none
\begin{question}{46}
국토의 계획 및 이용에 관한 법령상 기반시설에 관한 설명으로 옳지 않은 것은?
\begin{choices}
  \item 국가계획으로 설치ㆍ관리하는 광역시설은 그 광역시설의 설치ㆍ관리를 사업종목으로 하여 다른 법률에 따라 설립된 법인이 관리할 수 있다.
  \item 교통시설에 해당하는 자동차검사시설과 유통ㆍ공급시설에 해당하는 공동구는 기반시설에 속한다.
  \item 특별시장ㆍ광역시장ㆍ시장 또는 군수는 개발행위가 집중되어 해당 지역의 계획적 관리가 필요하다고 인정하는 지역을 기반시설부담구역으로 지정할 수 있다.
  \item 상업지역에서 지상ㆍ지하 등에 공공공지ㆍ열공급설비등을 설치하려고 할 경우 그 시설의 종류ㆍ명칭ㆍ위치ㆍ규모 등을 미리 도시관리계획으로 결정하여야 한다.
  \item 국가 또는 지방자치단체는 행정청이 아닌 자가 시행하는 도시계획시설사업에 드는 비용의 일부를 보조ㆍ융자할 수 있으며, 이때 해당지역의 기반시설이 인근지역에 비하여 부족한 지역을 우선 지원할 수 있다.
\end{choices}
\end{question}
\answer{4}
